\documentclass[11pt, letterpaper]{article}
\usepackage[margin=0.8in, top=0.6in, bottom=0.6in]{geometry}
\usepackage[utf8]{inputenc}
\usepackage{hyperref}

\hypersetup{
    colorlinks=true,
    linkcolor=blue,
    urlcolor=blue,
}

\setlength{\parskip}{0.6em}

\begin{document}

\thispagestyle{empty}

\begin{flushleft}
Yichao Jin, Ph.D.\\
School of Economic, Political and Policy Sciences\\
The University of Texas at Dallas\\
\href{mailto:Yichao.Jin@UTDallas.edu}{Yichao.Jin@UTDallas.edu}\\
\href{https://yichao2022.github.io}{https://yichao2022.github.io}
\end{flushleft}

\vspace{0.3em}

\begin{flushleft}
July 29, 2026
\end{flushleft}

\vspace{0.3em}

\begin{flushleft}
\textbf{Drs. Scott Halpern, Michael Harhay, and Ari Friedman}\\
Perelman School of Medicine\\
University of Pennsylvania
\end{flushleft}

\vspace{0.3em}

\textbf{Re: Postdoctoral Researcher in Health Services Research and Health Economics}

\vspace{0.3em}

Dear Drs. Halpern, Harhay, and Friedman,

I am writing to apply for the Postdoctoral Researcher position in Health Services Research and Health Economics. My research combines health economics, behavioral decision science, and causal inference to understand how healthcare policies and clinical interventions influence patient decision-making. I am excited by the opportunity to contribute to your collaborative research program, whose integration of health services research, behavioral science, and quantitative methodology closely aligns with my research interests.

My doctoral research examined behavioral barriers to preventive healthcare using discrete choice experiments and structural econometric models. Using data from 1,027 adults, I investigated how waiting time, institutional trust, and treatment attributes jointly shape vaccination decisions and identified behavioral mechanisms relevant to health policy and intervention design. Building on this work, I developed quantitative methods that integrate AI-derived behavioral information with empirically estimated choice models while preserving statistical validity and behavioral consistency. This work has been submitted to \textit{Value in Health} and reflects my broader interest in developing rigorous methods for evidence-based healthcare decision-making.

My methodological training includes health economics, causal inference, structural econometrics, and statistical programming in R and Python. Throughout my doctoral research, I independently designed empirical studies, managed complex datasets, developed reproducible analytical workflows, and led projects from study design through manuscript preparation. Working across economics, public policy, and health research has also strengthened my ability to collaborate with interdisciplinary teams and communicate quantitative methods to diverse audiences.

I am particularly attracted to this position because of its emphasis on rigorous causal inference using complex healthcare data to address clinically important questions. I look forward to contributing my background in behavioral economics and quantitative methods while expanding my expertise in electronic health record research, natural experiments, and cost-effectiveness analysis. I would be honored to contribute to your research programs and appreciate your consideration of my application.

\vspace{0.5em}

Sincerely,

\vspace{1em}

\textbf{Yichao Jin, Ph.D.}

\end{document}
